% =============================================================================
% test-abnt-unb-A.tex
% Teste visual — abnt-unb-A.cls (Modelo UnB — engenharia)
%
% Compilação (da raiz do repositório):
%   TEXINPUTS=./src:./assets/unb: latexmk -pdf -outdir=tests tests/test-abnt-unb-A.tex
%
% Checklist de verificação visual:
%   1  — Logo UnB centralizado no topo da capa
%   2  — "UNIVERSIDADE DE BRASÍLIA" em maiúsculas abaixo do logo
%   3  — Nome da faculdade abaixo da instituição
%   4  — Nome do departamento abaixo da faculdade
%   5  — Autor centralizado em maiúsculas
%   6  — Título LARGE, negrito, maiúsculas, centralizado
%   7  — Subtítulo em peso normal, precedido por ":"
%   8  — Local e ano no rodapé da capa
%   9  — Folha de rosto: logo menor (1.5cm), cabeçalho institucional em \small
%  10  — Folha de rosto: natureza do trabalho alinhada à direita
%  11  — Folha de aprovação funciona (herdada)
%  12  — Dedicatória funciona (herdada)
%  13  — Agradecimentos funciona (herdada)
%  14  — Epígrafe funciona (herdada)
%  15  — Resumo e abstract funcionam (via abnt-resumo)
%  16  — Sumário funciona (herdado)
%  17  — Compila com pdfLaTeX (font fallback TeX Gyre Termes)
% =============================================================================

\documentclass{abnt-unb-A}
\usepackage{abnt-resumo}

% =============================================================================
% METADADOS UnB
% =============================================================================
\faculdade{Faculdade UnB Gama}
\departamento{Departamento de Engenharia de Software}
\curso{Engenharia de Software}
\grau{Bacharel}
\tipotrabalho{Trabalho de Conclusão de Curso}

% =============================================================================
% METADADOS ABNT (herdados de abnt.cls)
% =============================================================================
\titulo{Sistema de monitoramento ambiental usando IoT}
\subtitulo{uma abordagem baseada em microcontroladores}
\autor{Fulana de Tal Silva}
\orientador{Prof.\ Dr.\ Beltrano Souza}
\ano{2025}

\natureza{Trabalho de Conclusão de Curso apresentado como requisito
parcial para obtenção do grau de Bacharel em Engenharia de Software
pela Universidade de Brasília.}
\programa{Engenharia de Software}
\area{Sistemas Embarcados}

% Banca examinadora
\dataaprovacao{20 de março de 2025}
\membrodabanca{Prof.\ Dr.\ Beltrano Souza}
  {Universidade de Brasília (Orientador)}
\membrodabanca{Prof.\ Dr.\ Sicrana Oliveira}
  {Universidade de Brasília}
\membrodabanca{Prof.\ Dr.\ Ciclano Pereira}
  {Universidade de Brasília}

% =============================================================================
% DOCUMENTO
% =============================================================================
\begin{document}

% --- Pré-textuais ---
\pretextual

\imprimircapa
\imprimirfolhaderosto
\imprimirfolhadeaprovacao

\imprimirdedicatoria{%
Aos meus pais, pelo apoio incondicional durante toda a graduação.}

\imprimiragradecimentos{%
Agradeço ao meu orientador, Prof.\ Dr.\ Beltrano Souza, pela
orientação dedicada ao longo deste trabalho.

Agradeço também aos colegas do Laboratório de Sistemas Embarcados
pela colaboração nos testes de campo.}

\imprimirepigrafe{%
\textit{``A ciência nunca resolve um problema sem criar pelo menos
outros dez.''} --- George Bernard Shaw}

\imprimirresumo{%
Este trabalho apresenta o desenvolvimento de um sistema de
monitoramento ambiental baseado em dispositivos IoT utilizando
microcontroladores de baixo custo.\ O sistema proposto realiza a
coleta contínua de dados de temperatura, umidade e qualidade do ar
em ambientes urbanos, transmitindo as informações para uma plataforma
em nuvem para análise e visualização em tempo real.\ A validação
experimental demonstrou que o sistema é capaz de operar de forma
autônoma por períodos prolongados com consumo energético reduzido,
atendendo aos requisitos de monitoramento ambiental definidos pela
legislação brasileira vigente.%
}{IoT; monitoramento ambiental; sistemas embarcados; microcontroladores;
qualidade do ar}

\imprimirabstract{%
This work presents the development of an environmental monitoring
system based on IoT devices using low-cost microcontrollers.\ The
proposed system continuously collects temperature, humidity, and air
quality data in urban environments, transmitting the information to a
cloud platform for real-time analysis and visualization.\ Experimental
validation demonstrated that the system can operate autonomously for
extended periods with reduced energy consumption, meeting the
environmental monitoring requirements defined by current Brazilian
legislation.%
}{IoT; environmental monitoring; embedded systems; microcontrollers;
air quality}

\imprimirsumario

% --- Textuais ---
\textual

\chapter{Introdução}

O monitoramento ambiental é uma atividade essencial para a gestão
sustentável dos recursos naturais e para a proteção da saúde humana.
Com o avanço da Internet das Coisas (IoT), tornou-se viável a
implantação de redes de sensores de baixo custo capazes de coletar
dados ambientais de forma contínua e distribuída.

\section{Objetivos}

\subsection{Objetivo geral}

Desenvolver um sistema de monitoramento ambiental baseado em
microcontroladores de baixo custo integrados a uma plataforma IoT.

\subsection{Objetivos específicos}

\begin{alineas}
  \item projetar o hardware de coleta de dados ambientais;
  \item implementar o firmware de comunicação sem fio;
  \item desenvolver a interface de visualização em nuvem;
  \item validar o sistema em ambiente real.
\end{alineas}

\chapter{Fundamentação teórica}

A fundamentação teórica abrange os conceitos de Internet das Coisas,
protocolos de comunicação sem fio e técnicas de monitoramento ambiental.

\section{Internet das Coisas}

A IoT refere-se à interconexão de dispositivos físicos por meio de
redes de comunicação, permitindo a troca de dados e o controle remoto
de sistemas.

% --- Pós-textuais ---
\postextual

\chapter*{Referências}
\addcontentsline{toc}{chapter}{REFERÊNCIAS}

{\singlespacing
Referências seriam geradas aqui via \texttt{\textbackslash imprimirreferencias}
com um arquivo \texttt{.bib}.}

\end{document}
